% Isolated research model after CP17. No determinant-nonvanishing claim.
\subsection{A rank-two bundle model for the remaining trace variation}

\begin{lemma}[The two toral bundles differ by an acyclic line bundle]
\label{lem:DE-A3-positive-two-bundles}
\status{Candidate; complete proof written, independent audit pending.}
Let $B$ be a mixed abelian surface over an algebraically closed
field $k$ of characteristic two and let
$H=M^{3d}$ be ample with positive ordinary index. Put
$T=B^{(2)}$, $L=H^{(2)}$ and $b=h^0(B,H)$.
There are rank-two bundles $E_0,E_1$ on $T$ with
\[
 F_{B/k,*}H=E_0\oplus E_1,\qquad
 E_1\simeq E_0\otimes P,
\]
where $P$ is the nontrivial two-torsion point of $\Pic^0(T)(k)$.
The Cartier pairing is orthogonal for this decomposition and
restricts to perfect alternating pairings with values in $L$.
In particular $\det E_i\simeq L$. Moreover
\[
 H^q(T,E_i)=0\ (q>0),\quad h^0(T,E_i)=b/2,\quad
 c_2(E_i)=b/2,\quad H^q(T,P)=0\ \text{for all }q.
\]
Both $E_i$ are globally generated.
\end{lemma}
\begin{proof}
The previously proved integral mixed decomposition gives
$\ker F_{B/k}\simeq\mu_2\times\alpha_2$.
Factor relative Frobenius as
\[
 B\xrightarrow{\rho}C=B/\mu_2\xrightarrow{\nu}T.
\]
The second map is an $\alpha_2$-torsor.
Since $3$ is odd, positive ordinary index of $M^{3d}$ implies
that $M^d$ descends along $\rho$, by the actual theta-group
linearization already proved. Thus $H=\rho^*J$ with $J$ a
cube of an ample line bundle. The torsor algebra is
$\rho_*\mathcal O_B=\mathcal O_C\oplus Q$ for the nontrivial
two-torsion point $Q\in\Pic^0(C)(k)$, and therefore
\[
 E_0=\nu_*J,\qquad E_1=\nu_*(J\otimes Q).
\]
The dual of $\nu$ has kernel $\alpha_2$. Consequently
$\nu^*: \Pic^0(T)(k)\to\Pic^0(C)(k)$ is bijective.
The unique preimage $P$ of $Q$ is nontrivial and has order two:
$\nu^*(P^2)=\mathcal O_C$, and the kernel has only the identity
as a geometric point. Projection formula gives $E_1=E_0P$.

Finite Frobenius duality, or its elementary local two-basis
calculation, makes $C(fg\omega)$ a perfect pairing on $F_*H$
with values in $L$. Its diagonal is zero because
$C(\omega)=0$ on a mixed surface.
The toral operator is self-adjoint for this pairing:
$C(e(fg)\omega)=0$ gives $\beta(ef,g)=\beta(f,eg)$.
Its different eigenspaces are therefore orthogonal; each
restriction is perfect alternating of rank two. This proves
the determinant assertions.

Ample abelian-variety cohomology vanishing on $C$ gives the
asserted higher vanishing for $E_i$. Riemann--Roch and the
degree-two map $\rho$ give $h^0(E_i)=\chi(J)=b/2$.
On an abelian surface the tangent bundle is trivial, so
Riemann--Roch for a rank-two bundle gives
$\chi(E_i)=c_1(E_i)^2/2-c_2(E_i)$. Since $c_1(E_i)=c_1(L)$
and $L^2/2=b$, this proves $c_2(E_i)=b/2$.
The acyclicity of $P$ is
\cite[Lemma~7.19 and its full proof, p.~105]
{EdixhovenVanDerGeerMoonenAV}.
Finally $J$ and $JQ$ are cubes of ample line bundles, after
choosing a cube root of $Q$ in $\Pic^0(C)(k)$, and so are
very ample. They separate every length-two subscheme,
including every fibre of the finite flat map $\nu$.
Thus the global evaluations of $\nu_*J$ and $\nu_*(JQ)$
are surjective at every point.
\end{proof}

\begin{lemma}[An acyclic summand and a determinant curve]
\label{lem:DE-A3-positive-determinant-curve}
\status{Candidate; complete proof written, independent audit pending.}
Keep the preceding notation. A general pair
$s_0\in H^0(E_0)$, $s_1\in H^0(E_1)$ has disjoint finite
local-complete-intersection zero schemes $Z_0,Z_1$, each
of length $b/2$.
Normalize the isomorphism $E_1=E_0P$ to identify its alternating
pairing with that on $E_0$, using $P^2=\mathcal O_T$.
Then
\[
 K=\ker\bigl(E_0\oplus E_1
       \xrightarrow{\beta(s_0,-)+\beta(s_1,-)}L\bigr)
\]
is locally free of rank three. It contains the line subbundle
$P\to K$, $a\mapsto(as_1,as_0)$.
The quotient $Q_s=K/P$ has rank two, determinant $LP$, and
the same cohomology as $K$.

Put $w=s_0\wedge s_1\in H^0(LP)$. For general pairs $w\ne0$.
The two sections $(s_0,0)$ and $(0,s_1)$ induce an exact sequence
\[
 0\longrightarrow\mathcal O_T^2\longrightarrow Q_s
      \longrightarrow i_*R_s\longrightarrow0,
 \qquad i:D_s=(w=0)\hookrightarrow T,
                                                        \tag{DE-TB1}
\]
where $R_s$ is an invertible sheaf on the Cartier divisor $D_s$.
The arithmetic genus of $D_s$ is $b+1$, and
$\chi(D_s,R_s)=0$, so $\deg R_s=b$.
No reducedness of the zero schemes or of $D_s$ is assumed here.

If $s=s_0+s_1$ is a sharp source in the preceding variation
lemma, then $h^0(Q_s)=h^1(Q_s)=2$ and $h^2(Q_s)=0$.
\end{lemma}
\begin{proof}
Global generation gives the usual zero incidence over $T$.
For a general section its zero scheme has dimension zero;
the two local equations are a regular sequence on the smooth
surface and its length is $c_2(E_i)=b/2$.
Requiring simultaneous vanishing of $s_0,s_1$ imposes four
independent conditions over a two-dimensional base, so the
general two zero schemes are disjoint. The contraction map
onto $L$ is then surjective. The displayed map from $P$ lies
in its kernel by symmetry of the two cross terms in
characteristic two, and is a subbundle because $s_0,s_1$
never vanish together. This proves local freeness of $Q_s$.
Determinants give $\det Q_s=LP$; acyclicity of $P$ gives the
cohomology assertion.

At a fixed point, the two section evaluations can be chosen
independently, so $w$ is not identically zero for a general
pair. The determinant of the two induced sections of $Q_s$
is $w$, with the compatible determinant identifications.
One may check this in a frame in which
$s_0=(1,0)$ and $s_1=(f,w)$. If the chosen frame $e$ of $P$
satisfies $e^2\mapsto u\in\mathcal O_T^\times$ under the
fixed trivialization of $P^2$, the two columns in $Q_s$ become
$(1,0)$ and $(u f^2,w)$. It is not necessary, and need not
be possible etale locally, to make $u$ a square.
Their determinant is still $w$. The chart with $s_1$ nonzero treats
the other points. Along $D_s$ the two columns always have
rank one, since the original sections do not vanish together.
Thus their cokernel is a line bundle on $D_s$, proving
(DE-TB1), even if the divisor is nonreduced.

Adjunction gives $p_a(D_s)=1+(LP)^2/2=b+1$.
The contraction sequence gives
$\chi(K)=\chi(E_0)+\chi(E_1)-\chi(L)=0$.
Then $\chi(Q_s)=0$ and (DE-TB1) gives $\chi(R_s)=0$.
Riemann--Roch for a line bundle on a projective Cartier
curve gives its asserted degree.
Finally sharpness says that global Cartier contraction has
kernel and cokernel of dimension two. Higher cohomology
of $E_i,L$ vanishes; the long exact sequence therefore gives
$h^0(K)=h^1(K)=2$, $h^2(K)=0$. The same holds for $Q_s$.
\end{proof}

\begin{lemma}[The actual torsor class in the contraction boundary]
\label{lem:DE-A3-positive-torsor-boundary}
\status{Candidate; complete proof written, independent audit pending.}
For an etale double cover $\pi':\widetilde T\to T$, let
$\alpha\in H^1(T,\mathcal O_T)$ be the extension class of
$0\to\mathcal O_T\to\pi'_*\mathcal O_{\widetilde T}
\xrightarrow{\operatorname{Tr}}\mathcal O_T\to0$.
Choose a trace lift $q_i$ of $s_i$ in
$H^0(T,E_i\otimes\pi'_*\mathcal O_{\widetilde T})$.
The contraction $\beta(s_i,q_i)$ descends to a section
$v_i\in H^0(T,L I_{Z_i})$. Its boundary in
\[
 0\longrightarrow\mathcal O_T\xrightarrow{s_i}E_i
       \xrightarrow{\beta(s_i,-)}L I_{Z_i}
       \longrightarrow0
\]
is exactly $\alpha$.
\end{lemma}
\begin{proof}
Trace commutes with contraction, and
$\operatorname{Tr}\beta(s_i,q_i)=\beta(s_i,s_i)=0$.
Thus contraction descends. Its vanishing on $Z_i$ is local.
Choose local trace-one generators $z_a$ of the torsor algebra,
with $z_b-z_a=c_{ab}$ representing $\alpha$.
Write $q_i=z_a s_i+r_a$ in each local trivialization.
Then $\beta(s_i,q_i)=\beta(s_i,r_a)$, and
$r_b-r_a=c_{ab}s_i$ (signs agree in characteristic two).
These are precisely the Cech representatives of the
connecting homomorphism. The resulting class is $\alpha$.
Existence of the trace lift follows from $H^1(T,E_i)=0$.
\end{proof}

\begin{remark}[The remaining step]
\label{rem:DE-A3-positive-bundle-scope}
\status{Candidate scope statement; no residue nonvanishing assertion.}
These statements give actual bundles and a Cartier curve on
which a residue calculation could be made. No formula has yet
identified the quadratic variation determinant with a
nonzero residue pairing on $(D_s,R_s)$. In particular neither
acyclicity of $P$ nor the degree of $R_s$ proves sharpness.
The exact outstanding assertion remains exclusion of a common
kernel or common image for the trace-direction matrix space.
\end{remark}
